\documentclass[10pt,a4paper]{amsart}
\usepackage{amsmath,amssymb,amsthm,graphicx}
\usepackage{verbatim}
\newenvironment{reference}{\comment}{\endcomment}
%\newenvironment{reference}{}{}
\newenvironment{slogan}{\comment}{\endcomment}
\newenvironment{history}{\comment}{\endcomment}
% For commutative diagrams we use Xy-pic
\usepackage[all]{xy}

% We use 2cell for 2-commutative diagrams.
\xyoption{2cell}
% Theorem environments.
%
\theoremstyle{plain}
\newtheorem{theorem}[subsection]{Theorem}
\newtheorem{proposition}[subsection]{Proposition}
\newtheorem{lemma}[subsection]{Lemma}

\theoremstyle{definition}
\newtheorem{definition}[subsection]{Definition}
\newtheorem{example}[subsection]{Example}
\newtheorem{exercise}[subsection]{Exercise}
\newtheorem{situation}[subsection]{Situation}

\theoremstyle{remark}
\newtheorem{remark}[subsection]{Remark}
\newtheorem{remarks}[subsection]{Remarks}

\numberwithin{equation}{subsection}

% Macros
%
\def\lim{\mathop{\mathrm{lim}}\nolimits}
\def\colim{\mathop{\mathrm{colim}}\nolimits}
\def\Spec{\mathop{\mathrm{Spec}}}
\def\Hom{\mathop{\mathrm{Hom}}\nolimits}
\def\Ext{\mathop{\mathrm{Ext}}\nolimits}
\def\SheafHom{\mathop{\mathcal{H}\!\mathit{om}}\nolimits}
\def\SheafExt{\mathop{\mathcal{E}\!\mathit{xt}}\nolimits}
\def\Sch{\mathit{Sch}}
\def\Mor{\mathop{\mathrm{Mor}}\nolimits}
\def\Ob{\mathop{\mathrm{Ob}}\nolimits}
\def\Sh{\mathop{\mathit{Sh}}\nolimits}
\def\NL{\mathop{N\!L}\nolimits}
\def\CH{\mathop{\mathrm{CH}}\nolimits}
\def\proetale{{pro\text{-}\acute{e}tale}}
\def\etale{{\acute{e}tale}}
\def\QCoh{\mathit{QCoh}}
\def\Ker{\mathop{\mathrm{Ker}}}
\def\Im{\mathop{\mathrm{Im}}}
\def\Coker{\mathop{\mathrm{Coker}}}
\def\Coim{\mathop{\mathrm{Coim}}}

% Boxtimes
%
\DeclareMathSymbol{\boxtimes}{\mathbin}{AMSa}{"02}

%
% Macros for moduli stacks/spaces
%
\def\QCohstack{\mathcal{QC}\!\mathit{oh}}
\def\Cohstack{\mathcal{C}\!\mathit{oh}}
\def\Spacesstack{\mathcal{S}\!\mathit{paces}}
\def\Quotfunctor{\mathrm{Quot}}
\def\Hilbfunctor{\mathrm{Hilb}}
\def\Curvesstack{\mathcal{C}\!\mathit{urves}}
\def\Polarizedstack{\mathcal{P}\!\mathit{olarized}}
\def\Complexesstack{\mathcal{C}\!\mathit{omplexes}}
% \Pic is the operator that assigns to X its picard group, usage \Pic(X)
% \Picardstack_{X/B} denotes the Picard stack of X over B
% \Picardfunctor_{X/B} denotes the Picard functor of X over B
\def\Pic{\mathop{\mathrm{Pic}}\nolimits}
\def\Picardstack{\mathcal{P}\!\mathit{ic}}
\def\Picardfunctor{\mathrm{Pic}}
\def\Deformationcategory{\mathcal{D}\!\mathit{ef}}
\usepackage{cleveref}
\crefname{originallemma}{original lemma}{original lemmas}
\Crefname{originallemma}{Original lemma}{Original lemmas}
\makeatletter
\newlabel{lemma-dualizing-finite-length}{{lemma-dualizing-finite-length}{source}{}{sourcepointer-1}{}}
\newlabel{lemma-dualizing-finite-length@cref}{{[originallemma][0][]lemma-dualizing-finite-length}{[1][0][]source}{}{}{}}
\makeatother
\crefname{originallemma}{original lemma}{original lemmas}
\Crefname{originallemma}{Original lemma}{Original lemmas}
\makeatletter
\newlabel{lemma-local-cohomology-ext}{{lemma-local-cohomology-ext}{source}{}{sourcepointer-2}{}}
\newlabel{lemma-local-cohomology-ext@cref}{{[originallemma][0][]lemma-local-cohomology-ext}{[1][0][]source}{}{}{}}
\makeatother
\crefname{originallemma}{original lemma}{original lemmas}
\Crefname{originallemma}{Original lemma}{Original lemmas}
\makeatletter
\newlabel{lemma-local-cohomology-noetherian}{{lemma-local-cohomology-noetherian}{source}{}{sourcepointer-3}{}}
\newlabel{lemma-local-cohomology-noetherian@cref}{{[originallemma][0][]lemma-local-cohomology-noetherian}{[1][0][]source}{}{}{}}
\makeatother
\crefname{originallemma}{original lemma}{original lemmas}
\Crefname{originallemma}{Original lemma}{Original lemmas}
\makeatletter
\newlabel{lemma-union-artinian}{{lemma-union-artinian}{source}{}{sourcepointer-4}{}}
\newlabel{lemma-union-artinian@cref}{{[originallemma][0][]lemma-union-artinian}{[1][0][]source}{}{}{}}
\makeatother
\title{Local cohomology of a normalized dualizing complex is the injective hull in degree zero}
\author{The Stacks Project Authors}
\date{stacks-project:dualizing.tex:local-snapshot-20261012; GNU FDL 1.2 or later}
\begin{document}
\maketitle
\noindent\textit{Editorial note.} All delivered mathematical text is the author's own, copied from the pinned Stacks Project chapter source sources/stacks-dualizing/pinned-original-source.tex; no statement or proof step was rewritten, reordered or removed.\par\medskip
\noindent\textit{Editorial note.} Preamble provenance: the Stacks Project keeps one single shared preamble for every chapter; this card takes it directly from the staged sources/stacks-dualizing/preamble.tex as three verbatim spans (lines 8-12, 14-18 and 154-217), so no derived file was written and nothing was added or redefined. The wrapper class and the base packages are supplied by the assembler.\par\medskip
\noindent\textit{Editorial note.} Bibliography provenance: neither the counted statement nor its proof contains a citation (this lane rescans the whole span before writing the card), so the card carries no bibliography and no entry was reconstructed.\par\medskip
\noindent\textit{Editorial note.} Selection provenance: this card was opened under the file-based control contract (CONTROL.md), with the mechanical pre-card gate run on its statement and proof interval before the card was written.\par\medskip
\noindent\textit{Allowed standard prerequisites (editorial).}\begin{itemize}
\item The identification of local cohomology with support in the closed point with the local cohomology of the maximal ideal. Invoked at the start of the counted proof (source line 3651). Rewrites the complex whose cohomology is computed.
\item The computation of local cohomology as a homotopy colimit of internal homs out of the quotients by the powers of the ideal. Invoked inside the counted proof (source lines 3656-3657). Exhibits the complex as a homotopy colimit of computable pieces.
\item The identification of the internal hom out of a finite length quotient with the corresponding truncation of the injective hull. Invoked inside the counted proof (source lines 3660-3664). Replaces each piece of the homotopy colimit by a truncation of the hull, compatibly with the transition maps.
\item The description of the hull as the union of its truncations. Invoked inside the counted proof (source lines 3666-3667). Identifies the colimit with the hull and shows the other cohomology modules vanish.
\end{itemize}\medskip
\ifdefined\hypertarget\hypertarget{sourcepointer-1}{}\fi
\noindent\textit{Source reference.} Original lemma lemma-dualizing-finite-length: sources/stacks-dualizing/pinned-original-source.tex, source lines 3204--3204. Navigation pointer only: the counted text refers to the Stacks item lemma-dualizing-finite-length, which is an existing labeled result of this same chapter and is not part of this card. The pointer carries the target's own label as its number and its definition line as its span; its content is not reproduced here.\par\medskip
\ifdefined\hypertarget\hypertarget{sourcepointer-2}{}\fi
\noindent\textit{Source reference.} Original lemma lemma-local-cohomology-ext: sources/stacks-dualizing/pinned-original-source.tex, source lines 1308--1308. Navigation pointer only: the counted text refers to the Stacks item lemma-local-cohomology-ext, which is an existing labeled result of this same chapter and is not part of this card. The pointer carries the target's own label as its number and its definition line as its span; its content is not reproduced here.\par\medskip
\ifdefined\hypertarget\hypertarget{sourcepointer-3}{}\fi
\noindent\textit{Source reference.} Original lemma lemma-local-cohomology-noetherian: sources/stacks-dualizing/pinned-original-source.tex, source lines 1789--1789. Navigation pointer only: the counted text refers to the Stacks item lemma-local-cohomology-noetherian, which is an existing labeled result of this same chapter and is not part of this card. The pointer carries the target's own label as its number and its definition line as its span; its content is not reproduced here.\par\medskip
\ifdefined\hypertarget\hypertarget{sourcepointer-4}{}\fi
\noindent\textit{Source reference.} Original lemma lemma-union-artinian: sources/stacks-dualizing/pinned-original-source.tex, source lines 1042--1042. Navigation pointer only: the counted text refers to the Stacks item lemma-union-artinian, which is an existing labeled result of this same chapter and is not part of this card. The pointer carries the target's own label as its number and its definition line as its span; its content is not reproduced here.\par\medskip
\begin{lemma}
\label{lemma-local-cohomology-of-dualizing}
Let $(A, \mathfrak m, \kappa)$ be a Noetherian local ring.
Let $\omega_A^\bullet$ be a normalized dualizing complex.
Let $Z = V(\mathfrak m) \subset \Spec(A)$.
Then $E = R^0\Gamma_Z(\omega_A^\bullet)$ is an injective hull of
$\kappa$ and $R\Gamma_Z(\omega_A^\bullet) = E[0]$.
\end{lemma}

\begin{proof}
By Lemma \ref{lemma-local-cohomology-noetherian} we have
$R\Gamma_{\mathfrak m} = R\Gamma_Z$. Thus
$$
R\Gamma_Z(\omega_A^\bullet) =
R\Gamma_{\mathfrak m}(\omega_A^\bullet) =
\text{hocolim}\ R\Hom_A(A/\mathfrak m^n, \omega_A^\bullet)
$$
by Lemma \ref{lemma-local-cohomology-ext}. Let $E'$ be an injective
hull of the residue field.
By Lemma \ref{lemma-dualizing-finite-length}
we can find isomorphisms
$$
R\Hom_A(A/\mathfrak m^n, \omega_A^\bullet) \cong \Hom_A(A/\mathfrak m^n, E')[0]
$$
compatible with transition maps. Since
$E' = \bigcup E'[\mathfrak m^n] = \colim \Hom_A(A/\mathfrak m^n, E')$
by Lemma \ref{lemma-union-artinian}
we conclude that $E \cong E'$ and that all other cohomology
groups of the complex $R\Gamma_Z(\omega_A^\bullet)$ are zero.
\end{proof}

\end{document}
